% =====================================================================
%  Proposal-scale adaptation for the streaming MCMC sampler
%  Three candidate methods, motivation, and recommendation.
%  Build:  pdflatex proposal_adaptation_methods.tex   (run twice for refs)
% =====================================================================
\documentclass[11pt]{article}

\usepackage[margin=1in]{geometry}
\usepackage{amsmath,amssymb}
\usepackage{booktabs}
\usepackage{enumitem}
\usepackage{xcolor}
\usepackage[hidelinks]{hyperref}

\newcommand{\code}[1]{\texttt{\small #1}}
\newcommand{\param}[1]{\textsf{#1}}
\newcommand{\pc}{\code{pertcoeff}}

\title{\bfseries Proposal-scale adaptation for the streaming MCMC sampler\\[3pt]
\large Three candidate methods, with a recommendation}
\author{OpenHydroTwin --- \code{DTStreamingMCMC}}
\date{\today}

\begin{document}
\maketitle

\begin{abstract}
The streaming sampler currently proposes all parameters jointly with a
random walk whose per-parameter step widths are scaled by a \emph{single
global factor} adapted toward a target acceptance rate. On an
ill-conditioned posterior --- one sharp direction alongside broad,
weakly-identified ones --- this single knob cannot serve both: the global
scale collapses chasing acceptance in the sharp direction, freezing the
broad directions, so chains stop mixing and never plateau. This note
defines the failure precisely, then lays out three standard remedies ---
(A) componentwise Metropolis-within-Gibbs, (B) diagonal adaptive scaling,
and (C) adaptive-covariance --- with their proposal mechanics, how each
sets the per-parameter (or per-direction) scale, cost, handling of the
model's log-normal parameters, and trade-offs. The recommendation, and
what is implemented (as a config-toggled alternative), is (C) with the
posterior covariance \emph{accumulated across all cycles}: since the
correlation structure is model-determined and stationary, accumulation
gives a large effective sample size that removes the full covariance's only
weakness, with (B) as its diagonal special case.
\end{abstract}

% =====================================================================
\section*{Notation}
% =====================================================================

\begin{center}
\begin{tabular}{@{}ll@{}}
\toprule
Symbol & Meaning\\
\midrule
$\psi=(\psi_1,\dots,\psi_d)$ & parameter vector being sampled\\
$d$ & number of parameters ($=n_p$; here $d=9$)\\
$\theta$ & the physical parameters (subset of $\psi$; the rest are error SDs)\\
$\pi(\psi)$ & target posterior; $\log\pi$ its log\\
$x,\,X$ & current / proposed chain position\\
$q(\cdot\mid\cdot)$ & proposal density\\
$a(x\to X)$ & Metropolis--Hastings acceptance probability, Eq.~\eqref{eq:mh}\\
$q(x\mid X)/q(X\mid x)$ & Hastings correction (unity for a symmetric proposal)\\
$z,\,z_i$ & standard-normal draws, $z\sim\mathcal N(0,I_d)$\\
$s_i$ & per-coordinate proposal scale (\code{pertcoeff}$_i$)\\
$\kappa$ & single global proposal scalar; $\kappa_0=2.38/\sqrt d$\\
$\eta$ & Robbins--Monro adaptation step size\\
$\hat a$ & observed joint acceptance rate over a block\\
$\hat a_i$ & observed per-coordinate acceptance rate (Method A)\\
$a^\star$ & target acceptance; $a^\star\!\approx\!0.234$ (joint), $a^\star_{1\mathrm D}\!\approx\!0.44$ (1-D)\\
$\hat\sigma_i$ & posterior marginal standard deviation of parameter $i$\\
$\hat\Sigma$ & posterior covariance (accumulated across cycles here)\\
$\Sigma_k$ & cycle $k$'s mean-centered covariance, Eq.~\eqref{eq:accum}\\
$L$ & lower-triangular Cholesky factor, $\hat\Sigma=LL^{\!\top}$\\
$\varepsilon I_d$ & ridge added for positive-definiteness; $I_d$ the $d\times d$ identity\\
$\phi_k^{(s)}$ & cycle $k$, draw $s$, in \emph{proposal space} (log for log-normal coords)\\
$\bar\phi_k$ & cycle $k$'s own sample mean (used to mean-center)\\
$n_k$ & number of pooled draws in cycle $k$\\
$W$ & accumulated sample weight in the running covariance\\
$\operatorname{diag}(\cdot)$, $\operatorname{chol}(\cdot)$ & diagonal matrix from a vector; Cholesky factor\\
$\mathcal N(\mu,\Sigma)$ & multivariate normal with mean $\mu$, covariance $\Sigma$\\
\bottomrule
\end{tabular}
\end{center}

% =====================================================================
\section{Setting and the observed failure}
% =====================================================================

Let $\psi=(\psi_1,\dots,\psi_d)\in\mathbb R^{d}$ be the parameter vector
($d=n_p$; here $d=9$) with target $\pi(\psi)\propto\exp[\log\pi(\psi)]$,
the kernel-weighted posterior. A chain at position $x$ proposes $X$ and
accepts with the Metropolis--Hastings probability
\begin{equation}
  a(x\to X) = \min\!\Big\{1,\ \frac{\pi(X)}{\pi(x)}\,
  \frac{q(x\mid X)}{q(X\mid x)}\Big\},
  \label{eq:mh}
\end{equation}
where $q$ is the proposal density and the ratio $q(x\mid X)/q(X\mid x)$ is
the Hastings correction (unity for a symmetric proposal).

\paragraph{Current proposal (\code{proposeFrom}).} Per coordinate, with
scale $s_i=\pc_i$ and $z_i\sim\mathcal N(0,1)$:
\begin{equation}
  \text{additive (normal/uniform prior):}\ \ X_i = x_i + s_i z_i,
  \qquad
  \text{multiplicative (log-normal prior):}\ \ X_i = x_i\,e^{\,s_i z_i}.
  \label{eq:current}
\end{equation}
The multiplicative form is an additive random walk in $\log\psi_i$; its
Hastings correction is $\prod_i (X_i/x_i)$ over the log-normal
coordinates, which the sampler already carries explicitly. All $s_i$ are
adapted by \emph{one common multiplier} toward a target acceptance rate
$a^\star$ (\code{adaptProposalScale}), frozen once a quorum forms.

\paragraph{The failure.} Diagnosed from \code{debug.log}: as the data
window slid, the posterior became ill-conditioned. The global scale was
driven to $s\!\approx\!1.5\times10^{-4}$ yet acceptance stayed at
$2$--$16\%$ and oscillated; chains crawled, failed the plateau motion
guard, and certification collapsed. The single scalar cannot match a
sharp direction (needs tiny steps) and a broad direction (needs large
steps) simultaneously --- shrinking it to accept in the former freezes
the latter. All three methods below give the proposal \emph{direction- or
coordinate-specific} scale so the two regimes decouple.

\paragraph{What the sampler already computes.} At each cycle
\code{computeSummary} produces, from the pooled samples, the
per-parameter standard deviations $\hat\sigma_i$ and the full
$d\times d$ posterior covariance $\hat\Sigma$. These are the natural
inputs for methods (B) and (C), and (like \pc) can be warm-started from
the previous cycle.

% =====================================================================
\section{Method A --- Componentwise Metropolis-within-Gibbs}
% =====================================================================

\paragraph{Idea.} Abandon the joint proposal: update one coordinate at a
time, each with its own accept/reject, so each parameter has a genuine
per-parameter acceptance rate to adapt against.

\paragraph{Mechanics.} A sweep visits coordinates $i=1,\dots,d$ (fixed or
random order). For coordinate $i$, propose $X_i$ from
Eq.~\eqref{eq:current} holding $\psi_{-i}$ fixed, and accept/reject via
Eq.~\eqref{eq:mh} using the full posterior (each proposal still requires a
forward solve). Track the per-coordinate acceptance $\hat a_i$ over a
window.

\paragraph{Setting the scale.} Each $s_i$ is adapted independently by
Robbins--Monro toward the one-dimensional optimum,
\begin{equation}
  \log s_i \leftarrow \log s_i + \eta\,(\hat a_i - a^\star_{1\mathrm D}),
  \qquad a^\star_{1\mathrm D}\approx 0.44,
\end{equation}
(the optimal acceptance for a single-coordinate random walk;
\cite{gelman1996}). Because each coordinate has its own $\hat a_i$, sharp
and broad parameters converge to very different $s_i$ automatically.

\paragraph{Cost \& fit here.} Clean and robust, but a full sweep now costs
$d$ forward solves instead of one --- a $\sim\!9\times$ increase in the
dominant expense. Under a fixed wall-clock budget that cuts the effective
sweep count by $\sim\!9\times$, which is severe for this application.
Correlated directions are also handled only indirectly (one coordinate at
a time mixes slowly along a correlated ridge).

% =====================================================================
\section{Method B --- Diagonal adaptive scaling (recommended first step)}
% =====================================================================

\paragraph{Idea.} Keep the joint proposal (one solve per step) but set each
coordinate's step width proportional to that parameter's \emph{posterior
marginal spread}. No per-parameter acceptance is needed: the geometry
supplies the per-coordinate scale directly.

\paragraph{Mechanics.} With a single global scalar $\kappa>0$ and the
per-parameter posterior standard deviation $\hat\sigma_i$,
\begin{equation}
  s_i = \kappa\,\hat\sigma_i,
  \qquad\text{i.e.}\qquad
  X = x + \kappa\,\operatorname{diag}(\hat\sigma_1,\dots,\hat\sigma_d)\,z,
  \label{eq:diag}
\end{equation}
so a sharp direction ($\hat\sigma_i$ small) automatically receives small
steps and a broad, weakly-identified direction ($\hat\sigma_i$ large)
receives large steps. Only the \emph{single} scalar $\kappa$ is then
adapted toward the high-dimensional optimum,
\begin{equation}
  \log\kappa \leftarrow \log\kappa + \eta\,(\hat a - a^\star),
  \qquad a^\star\approx 0.234,\quad
  \kappa_0 = \frac{2.38}{\sqrt d},
  \label{eq:kappa}
\end{equation}
using the ordinary joint acceptance $\hat a$ (\cite{roberts2001}). Because
the \emph{shape} is fixed by $\hat\sigma$, $\kappa$ can no longer freeze
the broad directions while chasing the sharp one --- the exact pathology
that broke the current sampler.

\paragraph{This is a symmetric proposal} in the additive coordinates
(Gaussian random walk), so no Hastings correction there; for log-normal
coordinates the width $\hat\sigma_i$ is measured in \emph{log space}
(std.\ dev.\ of $\log\psi_i$) and the multiplicative form of
Eq.~\eqref{eq:current} is retained, leaving the existing Jacobian
correction unchanged.

\paragraph{Where $\hat\sigma$ comes from.} The carried posterior spread
from the previous cycle (warm-started, exactly as \pc\ is today), refreshed
mid-cycle from the retained samples once a pool forms. On cycle~1 / before
a pool exists, fall back to a prior-range width per coordinate.

\paragraph{Cost \& fit here.} One solve per step (no cost increase); uses
$\hat\sigma_i$ already computed; a contained change to \code{proposeFrom}
plus a scalar adaptation. It fixes the diagonal ill-conditioning that
caused the failure. Its one limitation: it ignores \emph{correlations} ---
a stiff \emph{diagonal-of-a-rotated} direction (e.g.\ the
$n_{\text{eng}}$--$n_{\text{nat}}$ compensation) is not straightened by a
purely diagonal scaling.

% =====================================================================
\section{Method C --- Adaptive covariance, accumulated across cycles (implemented)}
% =====================================================================

\paragraph{Idea.} Precondition the joint proposal by the \emph{full}
posterior covariance, so both scales and correlations are matched --- the
gold standard for ill-conditioned, correlated targets.

\paragraph{Mechanics.} With $\hat\Sigma$ the (accumulated) posterior
covariance and $L$ its Cholesky factor ($\hat\Sigma = LL^{\!\top}$),
\begin{equation}
  X = x + \kappa\,L\,z,\qquad z\sim\mathcal N(0,I_d),
  \label{eq:full}
\end{equation}
which draws $X\sim\mathcal N(x,\ \kappa^2\hat\Sigma)$. Following adaptive
Metropolis \cite{haario2001}, $\kappa_0 = 2.38/\sqrt d$, with a ridge
$\varepsilon I_d$ guaranteeing positive-definiteness (and preventing a
collapsed non-identifiable direction from giving a singular, zero-width
proposal). The single scalar $\kappa$ is adapted per cycle toward
$a^\star\approx0.234$ as in Eq.~\eqref{eq:kappa}, so the proposal
\emph{shape} is set by $\hat\Sigma$ and its \emph{overall size} by
$\kappa$ --- the two decoupled. Since $\mathcal N(x,\kappa^2\hat\Sigma)$ is
symmetric, no Hastings correction is needed for additive coordinates; the
log-normal coordinates are handled by building $\hat\Sigma$ in the
transformed (log) space and proposing there, then mapping back (retaining
the multiplicative Jacobian).

\paragraph{Estimating $\hat\Sigma$: accumulate across all cycles, no
forgetting.} The posterior covariance separates into a \emph{shape}
(correlation structure --- which directions are stiff/sloppy, how
parameters compensate) and a \emph{location} (the mean). The shape is set
by the model's sensitivity structure $\partial(\text{outputs})/\partial
\theta$; because the \emph{model is fixed}, that structure is essentially
\textbf{stationary across cycles even as the parameter values drift}. The
location is what drifts. This licenses estimating $\hat\Sigma$ from
\emph{all} cycles rather than one:
\begin{equation}
  \Sigma_k = \frac{1}{n_k-1}\sum_{s}\big(\phi_k^{(s)}-\bar\phi_k\big)
             \big(\phi_k^{(s)}-\bar\phi_k\big)^{\!\top},
  \qquad
  \hat\Sigma \leftarrow \frac{W\hat\Sigma + (n_k-1)\Sigma_k}{W+(n_k-1)},
  \ \ W \mathrel{+}= n_k-1,
  \label{eq:accum}
\end{equation}
where $\phi_k^{(s)}$ is cycle $k$'s draw $s$ in \emph{proposal space} (log
for log-normal coordinates), $\bar\phi_k$ its \emph{own} cycle mean, and
$\Sigma_k$ is therefore \emph{mean-centered per cycle} --- removing the
drifting location so it cannot inflate the estimate. No forgetting factor
is applied: with a stationary shape, more cycles simply give a better
estimate, so the running average over the whole run yields a large
effective sample size and a well-conditioned $\hat\Sigma$ --- and
shrinkage toward the diagonal (needed only when samples are scarce)
becomes unnecessary; a small ridge suffices. Only the fast-timescale
overall scale is left to $\kappa$.

\paragraph{Cost \& fit here.} Still one solve per step; the extra work is a
$d\times d$ Cholesky per cycle and an $L z$ product per proposal ---
negligible at $d=9$. Caveats: strictly the covariance is a local
linearization, so the shape can shift somewhat if the parameters wander far
or the data window exercises different sensitivities --- but a proposal
preconditioner need only be \emph{approximately} right (MH accept/reject
corrects the rest), so the accumulated average structure is more than
adequate and its stability outweighs the small time-variation. The
earliest cold-start cycles have distorted spreads and are best excluded
from the accumulation (or down-weighted by ESS).

% =====================================================================
\section{Comparison}
% =====================================================================

\begin{center}
\begin{tabular}{@{}lccc@{}}
\toprule
 & \textbf{A. Componentwise} & \textbf{B. Diagonal} & \textbf{C. Full covariance}\\
\midrule
Proposal & one coord.\ at a time & joint, diag.\ scaled & joint, $\hat\Sigma$-scaled\\
Scale source & per-coord.\ acceptance $\hat a_i$ & marginal $\hat\sigma_i$ & covariance $\hat\Sigma$\\
Adapted quantity & $d$ scales $s_i$ & one scalar $\kappa$ & one scalar $\kappa$ (+ $\hat\Sigma$)\\
Target accept.\ & $\approx0.44$ (1-D) & $\approx0.234$ & $\approx0.234$\\
Solves per sweep & $d\times$ (expensive) & $1\times$ & $1\times$\\
Fixes diag.\ stiffness & yes & yes & yes\\
Fixes correlations & weakly & no & yes\\
Extra code & sweep restructure & small & small (Cholesky)\\
Hastings correction & unchanged & unchanged & unchanged (log-space)\\
\bottomrule
\end{tabular}
\end{center}

% =====================================================================
\section{Recommendation and implementation notes}
% =====================================================================

\textbf{Adopt Method~C with the all-cycles accumulated covariance of
Eq.~\eqref{eq:accum}.} Because the correlation \emph{shape} is
model-determined and stationary, accumulating it over the whole run
(mean-centered per cycle) gives a large effective sample size --- which
dissolves the one drawback of the full covariance (fragility on thin
pools) and makes shrinkage unnecessary. Method~B is then simply its
diagonal special case, and Method~A is ruled out by its $\sim\!9\times$
solve cost. This is implemented as a config-toggled alternative
(\param{mcmc\_proposal\_mode}: \code{"global"} $\to$ legacy,
\code{"covariance"} $\to$ this method) so it can be validated against the
current sampler before adoption, and, if effective, ported to the OHQ MCMC
module.

As implemented in \code{DTStreamingMCMC}:
\begin{itemize}[leftmargin=1.4em,nosep]
  \item \textbf{Accumulation:} \code{accumulateProposalCovariance} folds
  each cycle's pooled, mean-centered, proposal-space covariance into the
  running $\hat\Sigma$ (Eq.~\ref{eq:accum}); the accumulator persists across
  cycles (never reset in \code{initializeCycle}).
  \item \textbf{Bootstrap:} the proposal falls back to the legacy
  per-coordinate walk until the accumulated weight exceeds $\sim\!2d$, then
  switches to $\kappa L z$.
  \item \textbf{Log-normal parameters:} $\hat\Sigma$ is built in $\log\psi$
  for those coordinates; the multiplicative proposal is kept so the existing
  Jacobian correction is untouched.
  \item \textbf{Adaptation:} the single scalar $\kappa$ replaces the global
  width multiplier ($\kappa_0=2.38/\sqrt d$), adapted toward the target
  acceptance with the freeze-on-quorum retained and a $\kappa$
  floor/ceiling so it cannot run away.
  \item \textbf{Guardrails:} a ridge $\varepsilon I_d$ on $\hat\Sigma$ before
  the Cholesky; on a (rare) failed factorization the cycle falls back to the
  legacy proposal.
  \item \textbf{Follow-ups (not yet done):} persist $\hat\Sigma$ in the
  posterior snapshot so it survives restarts; optionally exclude the
  earliest cold-start cycles from the accumulation.
\end{itemize}

\begin{thebibliography}{9}
\bibitem{gelman1996} A.~Gelman, G.~O.~Roberts, W.~R.~Gilks (1996).
  Efficient Metropolis jumping rules. \emph{Bayesian Statistics 5}, 599--607.
\bibitem{roberts2001} G.~O.~Roberts, J.~S.~Rosenthal (2001).
  Optimal scaling for various Metropolis--Hastings algorithms.
  \emph{Statistical Science} 16(4), 351--367.
\bibitem{haario2001} H.~Haario, E.~Saksman, J.~Tamminen (2001).
  An adaptive Metropolis algorithm. \emph{Bernoulli} 7(2), 223--242.
\end{thebibliography}

\end{document}
